\documentclass[11pt]{article}
\usepackage[margin=1in]{geometry}
\usepackage{hyperref}
\usepackage{enumitem}
\usepackage{titlesec}
\usepackage{fancyhdr}
\usepackage{booktabs}

% Header/Footer
\pagestyle{fancy}
\fancyhf{}
\rhead{CS 4300/5300 Software Engineering}
\lhead{Sprint 0-2}
\rfoot{Page \thepage}

% Section formatting
\titleformat{\section}{\Large\bfseries}{\thesection}{1em}{}
\titleformat{\subsection}{\large\bfseries}{\thesubsection}{1em}{}

\begin{document}

\begin{center}
{\LARGE\textbf{CS 4300/5300 Software Engineering}}\\[0.5em]
{\Large Sprint 0-2: User Stories, Lo-Fi UI, and Storyboards}\\[1em]
\end{center}

\section{Purpose}

You have been assigned a customer. Now you turn their ideas into work your team can begin working on. This is the largest Sprint 0 deliverable, and the one your next four sprints depend on. A weak backlog here produces four sprints of confusion.

\section{Sprint 0-2 Objectives}

\begin{center}
\begin{tabular}{lc}
\toprule
\textbf{Requirement} & \textbf{Points} \\
\midrule
High-level user stories (at least 4) & 25 pts \\
Story decomposition (at least 6) & 25 pts \\
Lo-fi UI alignment & 15 pts \\
Storyboard clarity & 10 pts \\
Gherkin acceptance criteria & 10 pts \\
Sprint planning logic & 10 pts \\
Cohesion and professional quality & 5 pts \\
\midrule
\textbf{Total} & \textbf{100 pts} \\
\bottomrule
\end{tabular}
\end{center}

\section{Detailed Requirements}

\subsection{High-Level User Stories (25 pts)}
Write at least \textbf{four} high-level stories describing the core capabilities of your application. Each must use the Connextra format:

\begin{center}
\textit{As a} [user], \textit{I want to} [goal], \textit{So that} [reason/value].
\end{center}

Each story must be SMART---specific, measurable, achievable, relevant, and scoped to a meaningful user-visible capability. You will lose points for stories that omit the ``So that'' clause, state vague goals such as ``manage stuff,'' are too granular to be high-level, or are too broad to finish in the semester.

\subsection{Story Decomposition (25 pts)}
Break your high-level stories into at least \textbf{six} decomposed sized with gherkin scenarios to fit inside a single sprint. Each decomposed story must trace clearly back to one of your high-level stories, and must also use the Connextra format. Separate your high-level stories from your decomposed stories clearly in the submission.

\subsection{Lo-Fi UI Alignment (15 pts)}
Produce low-fidelity mockups covering your high-level stories. Show the major screens, the navigation between them, and the key user actions on each.

These should be sketches or wireframes. Polished high-fidelity mockups are not what is being asked for, and they cost you time you do not have.

\subsection{Storyboard Clarity (10 pts)}
For each high-level story, write a brief storyboard explaining how the user enters the workflow, what actions they take, what the system returns, and what successful completion looks like.

\subsection{Gherkin Acceptance Criteria (10 pts)}
Write acceptance criteria as Gherkin scenarios---\textbf{Given}, \textbf{When}, \textbf{Then}---on your decomposed stories in GitHub Projects. These scenarios become the tests you write in Sprint 1, so write them as statements about observable behavior. You will be demonstrating these scenarios in class throughout the semester.

\subsection{Sprint Planning Logic (10 pts)}
Set up a GitHub Projects board containing every story in a backlog, and distribute your decomposed stories across four sprint groupings. Your plan should deliver visible user value early, build functionality incrementally, and avoid concentrating unrealistic scope in any one sprint. Putting every story in Sprint 1 and leaving later sprints empty is not a plan.

\subsection{Cohesion and Professional Quality (5 pts)}
Combine all artifacts into one PDF. Sections should align with each other---the stories, the UI, the storyboards, and the sprint plan should all describe the same application.

\section{Submission Requirements}

\subsection*{On Canvas (1 person per team)}
\begin{itemize}
    \item Submit a single document containing all artifacts above
\end{itemize}

\subsection*{On GitHub Projects}
\begin{itemize}
    \item All stories entered as cards in the backlog
    \item Gherkin acceptance criteria recorded on decomposed stories
\end{itemize}

\subsection*{AI Documentation}
Document in your project README where and when you got help from AI, and include the transcript URL for each instance of support.

\end{document}
